% ============================================================
%  QUOTATION MARKS (corrected 2026-09-30): Fysisk Tidsskrift prints guillemets
%    pointing inward, »…« (checked at 250--600 dpi, scan PDF pp. 331, 333). They
%    are coded \guillemotright…\guillemotleft (LaTeX's \guillemotleft is «);
%    until 2026-09-30 they were coded the other way round and printed «…».
%    Open question: in the body text the marks may be SINGLE guillemets ›…‹
%    (the binarised scan does not settle it); the title on p. 331 is double.
%  H. M. Hansen, reply to Kroman  [TITLE NOT YET VERIFIED]
%  Fysisk Tidsskrift, 15. Aargang (1916-17), c. pp. 14-22.
%
%  DIPLOMATIC TRANSCRIPTION (Danish)
%  Status: COMPLETE. Printed pp. 18-29, all 12 pages, transcribed 2026-09-04
%  from the scan described below. No gaps, no duplicates; braces and guillemets
%  balanced (7/7); compiles with 0 errors, 0 missing characters, 0 overfull and
%  0 underfull boxes. English translation not begun.
%
%  The physicists' main answer to Kroman, and the fifth item in the debate --
%  the »tredie Afhandling« Kroman singles out in 1917: »Med Forfatteren af den
%  tredie Afhandling, Hr. Dr. H. M. Hansen, er jeg derimod erkendelsesteoretisk
%  set meget uenig.« Hansen's opening registers the pile-up of contributions:
%  »Den Interesse for Relativitetsspørgsmaalet, som flere Artikler i Fysisk
%  Tidsskrift i den senere Tid har givet Vidnesbyrd om...«
%  Hans Marius Hansen (1886-1956) was Niels Bohr's collaborator; Bohr had also
%  sat Kroman's filosofikum. See ../relativitetsprincipet/note.md.
%
% ------------------------------------------------------------
%  *** RIGHTS -- IN COPYRIGHT UNTIL 31 DECEMBER 2026 ***
% ------------------------------------------------------------
%  H. M. Hansen died 13 June 1956. Danish copyright is life + 70, so this text
%  is PUBLIC DOMAIN FROM 1 JANUARY 2027 and not before. Transcribing it now is
%  fine; publishing it is not.
%
%  Two things hold it back, and BOTH must be undone on release day:
%    1. \embargotrue in the preamble, which prints the notice on the first page.
%    2. this article's line in the repo's .gitignore, under EMBARGOED TEXTS.
%  Note that .gitignore does not cover files already tracked -- this file and
%  its PDF were, and had to be `git rm --cached`ed on 2026-09-04 before the
%  exclusion could bite.
%
% ------------------------------------------------------------
%  THE WITNESS
% ------------------------------------------------------------
%  Google Books `Tr4ZAAAAIAAJ` -- Fysisk tidsskrift, VOLUMES 13-15, the
%  University of California copy, digitised 17 January 2007, 886 pages, full
%  view. Kept as texts/kroman/scan-ft-13-15.pdf, untracked; shared with the
%  sibling articles. All twelve pages were inspected as images.
%
%  THIS SCAN LOSES HAIRLINES to its JBIG2 compression, and that matters more
%  here than anywhere else in the debate, because this is the one article by a
%  working physicist. Three losses were identified and are NOT transcribed as
%  defects, because they are the scan's and not the printing's:
%    p. 18  the rule of the fraction (v/c)^2 is simply gone at native 600 ppi.
%           Confirmed real by calibrating against the surviving rule in the
%           p. 20 footnote and the 1/2 rules on p. 27; set with \dfrac.
%    p. 19  the second d of »uladede«
%    p. 21  the ø slash throughout the footnote -- which is why »føre« there is
%           NOT logged as a printer's error
%  Check every thin rule at a second resolution before believing it absent.
%
% ------------------------------------------------------------
%  PAGE MAP -- VERIFIED 2026-09-04
% ------------------------------------------------------------
%      printed 18-29  =  PDF + 602   (PDF 620-631), uniform.
%
%  The same offset as Holst's »Tidsproblemet« at the front of the volume, and
%  confirmed at both ends: PDF 620 = printed 18 opens with the title block
%  »Relativitetsprincipet. / Af / H. M. Hansen.«, and printed p. 30 begins
%  »Fysik i Skolen«. NOTE the offset is NOT constant across vol. 15 -- it is
%  618 by p. 165 and 622 by p. 192 -- because unpaginated advertisement leaves
%  are bound in. Measure it per article; never carry one across.
%
%  Unlike three of its neighbours, this article does NOT begin partway down the
%  previous article's last page: it starts at the head of p. 18.
%
% ------------------------------------------------------------
%  CONVENTIONS -- SETTLED AGAINST THE IMAGES
% ------------------------------------------------------------
%  * Diplomatic. 1916 orthography exactly as printed; the pre-1948 aa kept,
%    nouns capitalised, nothing modernised.
%  * QUOTATION MARKS: guillemets »...«, 7 pairs, balanced.
%  * EMPHASIS is LETTERSPACING, set \emph{}; no true italic anywhere. Hansen's
%    compositor has his own rule, different again from Kroman's and Holst's:
%    he spaces a proper name ON FIRST SIGNIFICANT MENTION and leaves it plain
%    thereafter -- Kroman, Helge Holst, Lorentz-Fitzgerald, Poincaré, Rayleigh,
%    Brace, Trouton, Noble, Rankine, Einstein, Ritz, Wiechert, Schaefer,
%    Neumann, Kaufmann, Guye, Lavanchy, Ehrenfest. pp. 20, 22, 24 and 25 carry
%    no spaced names at all. Four ordinary words are spaced for stress: kan
%    (pp. 21, 28), sit (p. 26), Nøjagtighed (p. 28). Suffixes fall OUTSIDE the
%    run (\emph{Lorentz}', \emph{Kroman}'s) with one exception, \emph{Einsteins},
%    which is spaced whole.
%  * FOOTNOTES: *) and **), RESTARTING on each printed page -- ten notes, on
%    pp. 18 (two), 19, 20 (two), 21, 22, 23 and 27 (two). \thefootnote takes the
%    \ifcase form and \setcounter{footnote}{0} is issued at each page start that
%    needs it. A fixed *) is wrong here.
%  * MATHEMATICS: inline throughout. NO displays, and therefore no equation
%    numbers anywhere. Variables italic as printed; \dfrac for the two built-up
%    fractions, \tfrac for the case fractions of pp. 27-28; \gamma for the
%    acceleration; $43''$ for the Mercury perihelion arcseconds. A and B are
%    italic in the body but ROMAN in the p. 21 footnote, and are set as printed.
%  * The old Danish id-est sign ɔ: (p. 18) is built from a reflected c; see the
%    preamble.
%
% ------------------------------------------------------------
%  PRINTER'S DEFECTS -- transcribed AS PRINTED, logged at the site
% ------------------------------------------------------------
%    p. 28  »midlertid« for »imidlertid«
%    p. 28  »adskillgt« for »adskilligt«
%    p. 23, and the first note of p. 27: the footnote mark prints a bare * with
%           no closing parenthesis, against *) everywhere else
%  Recorded as readings rather than defects: the running head of p. 28 reads
%  »H. M Hansen:« with the point after the M dropped, and p. 19 carries the
%  gathering signature »2*«.
%
% ------------------------------------------------------------
%  HOW IT ENDS
% ------------------------------------------------------------
%  On p. 29, after a \medskip-sized break before the final two paragraphs. No
%  closing rule, no signature, no dateline; the rest of the page is blank.
%  There are no figures or tables anywhere in the article.
%
% ------------------------------------------------------------
%  BUILD
% ------------------------------------------------------------
%  pdflatex path (libertinus + libertinust1math), so the Makefile's default
%  pattern rule builds it. Sandbox compile-test per TRANSCRIPTION-PLAYBOOK.md §5.
% ============================================================

\documentclass[12pt,a4paper]{article}
\usepackage[utf8]{inputenc}
\usepackage[T1]{fontenc}
\usepackage{libertinus}
\usepackage{libertinust1math}
\usepackage{textalpha}   % Greek in text, typed directly
\usepackage{amsmath}     % this debate has real formulas; see CONVENTIONS

\usepackage{graphicx}    % for \reflectbox, below
% The old Danish id-est sign »ɔ:« (p. 18). T1 has no OPEN O, so it is built
% from a reflected c, as texts/brochner/det-religioese does.
\DeclareUnicodeCharacter{0254}{\reflectbox{c}}

% FOOTNOTES. The printing marks them *) and **), restarting on each printed
% page. A fixed mark will not do here, so the counter is reset at the page
% starts that need it (the resets are in the body, at the page markers).
\renewcommand{\thefootnote}{\ifcase\value{footnote}\or *)\or **)\else \arabic{footnote})\fi}
\usepackage[danish]{babel}
\usepackage[protrusion=true,expansion=true]{microtype}
\usepackage{setspace}
\usepackage[margin=1.2in]{geometry}
\usepackage{fancyhdr}
\setlength{\headheight}{15pt}
\usepackage{hyperref}

\hypersetup{
  pdftitle={H. M. Hansens svar til Kroman},
  pdfauthor={H. M. Hansen},
  hidelinks
}

\pagestyle{fancy}
\fancyhf{}
\fancyhead[LE,RO]{\thepage}
\fancyhead[RE]{\textit{Svar til Professor Kroman}}
\fancyhead[LO]{\textit{H. M. Hansen}}

\setstretch{1.3}


% ------------------------------------------------------------
%  EMBARGO SWITCH.  In copyright until 31 December 2026.
%  On 1 JANUARY 2027 change \embargotrue to \embargofalse below, delete this
%  article's two lines from the repo's .gitignore, and it publishes normally.
%  Nothing else needs changing.
% ------------------------------------------------------------
\newif\ifembargo
\embargotrue

\title{\textbf{\guillemotright Relativitetsprincipet\guillemotleft}}
\author{H.\ M.\ Hansen}
\date{\textit{Fysisk Tidsskrift} 15 (1916--17), pp.~18--29}

% ---- original-page markers (paginate.py) ----
\usepackage{marginnote}
\newcommand{\opage}[1]{\marginnote{\footnotesize\textit{[#1]}}}
\newcommand{\apage}[1]{\marginnote{\footnotesize\textit{[#1]}}}
\begin{document}
\maketitle
\ifembargo
\begin{center}\fbox{\parbox{0.86\textwidth}{\small\centering
\textbf{READING COPY --- NOT FOR CIRCULATION.}\\[2pt]
In copyright in Denmark until 31 December 2026 (H.\,M.\ Hansen died
13 June 1956; the term is life\,+\,70). Public domain from 1 January 2027.
Made for private study; please do not post or forward.}}\end{center}
\bigskip
\fi

% --- p. 18 ---
\opage{18} % The article begins at the HEAD of p. 18, under the title block
% »Relativitetsprincipet.« / Af / H. M. Hansen. -- which is set by \maketitle
% and is NOT repeated here. No epigraph, no section number before the first
% paragraph; the numbered sections 1.--8. run inline at paragraph openings.
%
% FOOTNOTES on this page: two, printed *) and **).
\setcounter{footnote}{0}%
Den Interesse for Relativitetsspørgsmaalet, som flere Artikler i
Fysisk Tidsskrift i den senere Tid har givet Vidnesbyrd om, tilskynder
mig til at knytte nogle Bemærkninger til disse Artikler. Det forekommer
mig nemlig, at baade Professor \emph{Kroman}'s\footnote{Fys. Tidsskr.
XIV, 1, 1915--16.} og Redaktør \emph{Helge Holst}'s Kritik i nogen Grad
bygger paa Misforstaaelser især af Lorentz' Opfattelse.

1.\ Da man havde forsøgt at forklare det negative Resultat af
Michelson's Interferensforsøg ved den \emph{Lorentz\,-\,Fitzgerald}ske
% The spaced run breaks over the line as »L o r e n t z - F i t z g e r a l d-«
% / »ske«: the hyphen is inside the letterspacing, the suffix -ske is not.
Sammentrækningshypotese (1892), indvendte \emph{Poincaré} (1900),
% There is a speck after »Hypotese« that the Google text layer reads as a
% comma. It sits above the baseline and is detached; read as dirt, not a comma.
at en saadan ad hoc fremsat Hypotese var en utilfredsstillende
Forklaring: det kunde jo blive nødvendigt at indføre nye Hypoteser,
hver Gang der fremkom nye Kendsgerninger. Dette skete nu meget snart.
Lord \emph{Rayleigh} sluttede i 1902, at dersom den Lorentzske
Sammentrækning fandt Sted, maatte Jordens Bevægelse gøre et
gennemsigtigt Legeme dobbeltbrydende. En saadan Dobbeltbrydning søgte
han forgæves at paavise, og senere har \emph{Brace} gentaget Forsøget
med en overordentlig følsom Kompensator, men ligeledes forgæves.
\emph{Trouton} og \emph{Noble} søgte i 1903 at paavise et
Drejningsmoment, som efter Elektronteorien maa virke paa en ladet
Pladekondensator, naar Pladerne danner en Vinkel med Jordens
Translationsretning: Kondensatoren var anbragt paa en tilstrækkelig
følsom Snovægt: den skulde altsaa dreje sig, naar den blev opladet, men
Resultatet var ogsaa negativt. Alle de her omtalte Virkninger skulde
ligesom Stribeforskydningen i Michelsonforsøget være Virkninger af
2.~Orden, proportionale med Kvadrat ɔ: $\left(\dfrac{v}{c}\right)^{2}$
% SCAN DEFECT, not a printing defect: the rule of this fraction is GONE from
% the scan -- absent at native 600 ppi in the 1-bit JBIG2 mask, so no amount of
% rendering recovers it. Calibrated against the same fraction in the second
% footnote of p. 20, where the rule and the radical's vinculum both survive,
% and against the built-up ½ fractions of p. 27, whose rules survive. The
% printing had the rule; the compression dropped it. Set with \dfrac.
paa Forholdet mellem Jordens Translationshastighed og
Lyshastigheden\footnote{En Ændring af Lysteorien for at forklare
Michelson-Forsøget, som foreslaaet af Prof. Kroman (S.~7), vilde ogsaa
rammes af Poincaré's Indvending, da den jo næppe vilde klare
Trouton-Noble-Forsøget.}.

% --- p. 19 ---
\opage{19} % FOOTNOTES on this page: one, printed *).
% Bottom right of the page carries the gathering signature »2*«; not text.
\setcounter{footnote}{0}%
2.\ Disse negative Resultater\footnote{Hertil er senere (1908) kommen et
forgæves Forsøg af \emph{Trouton} og \emph{Rankine} paa at paavise en
Indflydelse af Jordens Translation paa en Leders Ledningsevne, nemlig en
Forskel i Ledningsevnen i Retning af og tværs paa Translationen. En
saadan Forskel kunde ventes under Antagelse af Kontraktionshypotesen,
saa dette Forsøg er nær beslægtet med Rayleigh-Braces.} foranledigede
(1904) \emph{Lorentz} til at søge at ændre Grundlaget for
Elektronteorien saaledes, at den blev i Stand til at forklare i størst
mulig Almindelighed, at elektromagnetiske Fænomener i et vilkaarligt
System foregaar uafhængigt af Systemets Bevægelse (naar denne er uden
Acceleration). For at naa dette Resultat maatte i det væsentlige gøres
følgende nye Antagelser: 1) Selve Elektronerne, der som hvilende (i
Forhold til Æteren) antages kugleformede, trækker sig ved jævn
Bevægelse sammen i Bevægelsesretningen i det Forhold, som den gamle
Lorentz-Fitzgeraldske Hypotese forlangte. 2) Alle Kræfter mellem
uladede Partikler indbyrdes (altsaa f.~Eks. ogsaa elastiske Kræfter)
% »uladede«: the second d has failed almost completely in the scan -- two
% specks at the baseline where the sort should be. The rest of the word is
% clean and the sense is not in doubt; read as printed, no defect logged.
eller mellem ladede og uladede Partikler forandres ved en Translation
paa samme Maade som elektriske Kræfter i et elektrostatisk System.

Disse Antagelser viste sig i Stand til at forklare, at det ikke er
muligt at paavise en jævn Translation i Forhold til Æteren. Men det er,
som vi skal se, rigtignok en ganske anderledes vidtgaaende Hypotese end
den oprindelige Kontraktionshypotese, og det er en Misforstaaelse at
tro, at Antagelsen af denne sidste alene er Lorentz' Standpunkt, eller
at den alene skulde være i Stand til at forklare alle de ovenfor
omtalte Forsøg, der volder Vanskelighederne.

3.\ \emph{Einstein} greb i 1905 uden at kende Lorentz' Arbejde Sagen an
paa en anden Maade. Han var ligesom Lorentz kommen til den Opfattelse,
at det ikke var tilfældigt, at netop de hidtil anstillede Forsøg paa at
paavise Jordens Translation havde givet negative Resultater, men at det
var det rimeligste at antage, at hvad Aarsagen dertil end var, var den
af en saadan Art, at alle Forsøg i den Retning maatte strande. Han
opstillede da --- vi kan maaske sige foreløbig --- dette som Princip
for at drage Konsekvenserne deraf; dette Relativitetsprincip kan
f.~Eks. formuleres saaledes: Naturlovene er uafhængige af Iagttagerens
Bevægelsestilstand, naar denne er uden Acceleration. Der kan altsaa
ikke eksistere noget fra de andre udmærket Koordinatsystem, hvorved man
kan definere ab\-%
% --- p. 20 ---
\opage{20}% FOOTNOTES on this page: two, printed *) and **).
\setcounter{footnote}{0}%
solut Bevægelse, men alle Systemer, der bevæger sig i Forhold til
hverandre med jævn Hastighed, er nøjagtig lige gode til Beskrivelse af
et hvilket som helst Naturfænomen.

Hertil føjer Einstein nu det saakaldte Dogme om Lyshastighedens
Konstans. Dermed kan det siges at forholde sig paa følgende
Maade\footnote{Fys. Tidsskr. X, 262, 1911--12.}: Ved Fastsættelse af,
ved hvilke Transformationsligninger man under Antagelse af
Relativitetsprincipets Gyldighed gaar over fra et af disse
ligeberettigede Systemer til et andet, d.~v.~s. fra de Erfaringer, som
en Iagttager i det ene System gør, slutter til dem, en Iagttager i det
andet System vil gøre angaaende det samme Fænomen, maa vi gaa ud fra et
eller andet bekendt Fænomen. Hvilket, maa efter Relativitetsprincipet
være ligegyldigt; var det ikke det, d.~v.~s. fik man forskellige
Transformationsligninger ved at gaa ud fra forskellige Fænomener, var
jo netop Relativitetsprincipet ikke gyldigt. Men det Fænomen, vi
vælger, maa være tilstrækkelig simpelt og kendt med tilstrækkelig stor
Nøjagtighed, og derfor bliver der i Virkeligheden kun Tale om eet
Fænomen, Lysets Forplantning i det tomme Rum (eller Luft). Om den ved
vi netop fra Michelsons Forsøg med en overordentlig stor Nøjagtighed,
at den for os, der følger med, finder Sted med samme Hastighed $c$ i
alle Retninger. Relativitetsprincipet maa da forlange, at enhver
Iagttager i ethvert af de berettigede Systemer skal gøre den samme
Erfaring, at Lysets Hastighed i Vakuum ogsaa for ham er ens i alle
Retninger og lig $c$\footnote{Det er yderst faa fysiske Fænomener, vi
kender med en tilsvarende Nøjagtighed. Derfor kan der ikke hentes nogen
Indvending mod Relativitetsprincipet af den Omstændighed (Kroman
S.~22--23), at f.~Eks. Aberrationen eller Dopplers Princip faar Udtryk,
der afviger fra de sædvanlige med en Faktor
$\sqrt{1-\left(\dfrac{v}{c}\right)^{2}}$, da vi ikke tør paastaa, at
vore sædvanlige Udtryk er rigtige, og den eksperimentelle Prøve af
disse Fænomener ikke kan afgøre Tilstedeværelsen af en Faktor, hvis
Værdi ligger saa nær ved 1.}.

4.\ Hvad angaar Udledelsen af Transformationsligningerne paa det
saaledes skabte Grundlag kan jeg nøjes med at henvise til Professor
Kromans elegante Fremstilling. Transformationsligningerne er identiske
med dem, Lorentz Aaret før havde opstillet; alle Konsekvenserne
(specielt alle iagttagelige Konsekvenser) bliver derfor ogsaa
identiske, det er kun Fortolkningen der bliver forskellig i Over\-%
% --- p. 21 ---
\opage{21}% FOOTNOTES on this page: one, printed *). Note that A and B are set ROMAN in
% the footnote and ITALIC in the body; both followed as printed.
\setcounter{footnote}{0}%
ensstemmelse med de ovenfor omtalte forskellige Udgangspunkter.
(Adskillige af Konsekvenserne var imidlertid ikke bleven opdagede af
Lorentz, men kom først frem ved Einsteins Arbejde.)

\emph{Lorentz}' Fortolkning er følgende: Vi kan benytte Æteren til deri
at fastlægge et Koordinatsystem og saaledes definere absolut Hvile og
absolut Bevægelse i Forhold til dette Koordinatsystem. Vi kan maaske
ogsaa definere absolut Samtidighed for Begivenheder paa forskellige
Steder og ordne Tidspunkter i Tidsfølge paa entydig Maade, idet vi i
det mindste kan forestille os vilkaarlig store Hastigheder (altsaa paa
lignende Maade som i Redaktør Helge Holsts Artikel). Men de i
2.~omtalte Grundhypoteser fører til, at naar man meddeler et System en
jævn Translation i Forhold til Æteren, forkortes alle Længder i
Bevægelsesretningen, og alle Ure giver sig til at gaa langsommere (de
elastiske Kræfter i f.~Eks. Uroen forandrer sig jo ogsaa), saa at
ethvert System efter Lorentz faar sin \guillemotright effektive
Tid\guillemotleft . Disse Ændringer, der jo er udtrykt ved de
Lorentz-Einsteinske Transformationsformler, tilfredsstiller derfor
ogsaa den af Einstein først paapegede Reciprocitet: Dersom $A$ og $B$
bevæger sig i Forhold til hinanden med en vis relativ Hastighed, saa at
de bevæger sig med forskellig Hastighed i Forhold til Æteren, eller at
en af dem maaske er i Hvile i Forhold til denne, maa $A$ mene nøjagtig
det samme om $B$ (d.~v.~s. $B$'s Maalestokke og Ure), som $B$ mener om
$A$. Det er derfor udelukket, at nogen af dem eller nogen anden
\emph{kan} fastslaa, hvem af dem der muligvis hviler i Forhold til
Æteren, eller hvem der rigtigt kan afgøre, hvorvidt to Begivenheder paa
to forskellige Steder er samtidige eller ej\footnote{Lorentz' ogsaa af
Professor Kroman omtalte Diskussion mellem A og B (Das
Relativitätsprinzip, S.~22--23) skal kun oplyse dette Forhold. Hvis der
efter Lorentz bliver mindre heftig Strid mellem A og B end efter
Einstein, ligger det kun deri, at hvis de er Lorentztilhængere, kan de
enes om en fælles ideal Definition af Hvile og Samtidighed; de kan
aldrig afgøre, hvis Opfattelse der er den ideale, eller om nogen er
det. Men da A og B er Fysikere, vil de være for fornuftige til at føre
% The slash of ø is invisible here (»fore«), but it is equally invisible in
% »Afgørelse« three lines below and in »gøre« -- the scan drops that hairline
% throughout the footnote type. Absence of the slash is therefore not
% evidence; read »føre«, no defect logged.
denne Strid; der er jo egentlig intet at strides om og i hvert Fald
ingen Mulighed for en Afgørelse, hverken efter Lorentz' eller Einsteins
Opfattelse.}. Begreberne absolut Hvile, absolut Bevægelse og absolut
Samtidighed er efter Lorentz rent teoretiske eller ideale Begreber uden
nogen praktisk Betydning.

\emph{Einsteins} Standpunkt er simplere: Efter ham er der ingen Mening
i at tale om absolut Hvile, absolut Bevægelse eller absolut
% --- p. 22 ---
\opage{22} % FOOTNOTES on this page: one, printed *).
\setcounter{footnote}{0}%
Samtidighed. At man ikke kan forbinde nogen Mening med de to første
Begreber, naar man stiller sig paa et rent mekanisk Standpunkt og tror
paa Newtons Mekanik, er kendt nok, og det har man forlængst vænnet sig
til (Professor Kromans Bemærkninger herom (S.~18--19) omtales
nedenfor). I hvert Fald i Mekaniken mangler vi det fra de andre
udmærkede Koordinatsystemer, der kunde definere absolut Hvile og
Bevægelse. Men Antagelsen af en Æter kunde synes at give os et saadant
absolut Koordinatsystem, saa at man kunde vente gennem optiske og
elektrodynamiske Forsøg at blive i Stand til at skelne mellem Hvile og
Bevægelse; imidlertid har de ofte omtalte Forsøg vist os, at det heller
ikke er Tilfældet. I Stedet for Mekanikens Relativitetsprincip har
Einstein derfor maattet sætte sit eget, der paastaar, at hverken
Mekaniken eller nogen anden Gren af Fysiken giver os et saadant absolut
Koordinatsystem, og i saa Fald kan man efter hans Opfattelse
overhovedet ikke forbinde nogen Mening med at tale om absolut Hvile og
Bevægelse. Naar man drager Konsekvenserne af Relativitetsprincipet,
viser det sig naturligvis ganske som hos Lorentz, at hvert af de
ligeberettigede Systemer faar sin individuelle Opfattelse af
Samtidighed for Begivenheder paa forskellige Steder. Følgen deraf maa
for Einstein være, at man heller ikke kan forbinde Mening med at tale
om en absolut Tid eller absolut Samtidighed.

Enhver Iagttager kan ogsaa efter Einsteins Opfattelse direkte
konstatere, at i Systemer, der bevæger sig i Forhold til ham, er
Maalestokkene forkortede og gaar Urene langsommere og har Stedtid
ganske som efter Lorentz' Opfattelse\footnote{Vi forudsætter
naturligvis her som i det følgende, at Iagttagerne raader over
tilstrækkelig fine Maaleinstrumenter, saa at Fænomenernes lille
Størrelsesorden ingen Betydning faar. Det vil forhaabentlig ikke
forstyrre, at jeg for ikke at gøre Artiklen unødvendig bred mange
Steder har undladt at omtale forskellige tilsvarende
Forudsætninger.}. Der er altsaa ikke Tale om Forkortninger o.~s.~v.
blot som Formodninger i Iagttagerens Bevidsthed. Naar der efter
Einstein ikke er Mening i at tale om virkelige Forkortninger o.~s.~v.,
ligger det i, at der intet absolut System haves, hvori de
\guillemotright oprindelige\guillemotleft{} Længder og Urgange kan
fastslaas. Men det maa erindres, at ogsaa efter Lorentz' Opfattelse er
de virkelige Forkortninger o.~s.~v. i et vist System ikke de samme som
dem, en Iagttager, i Forhold til hvem Systemet bevæger sig,
% --- p. 23 ---
\opage{23} % FOOTNOTES on this page: one. Its mark is printed as a BARE ASTERISK: the
% closing paren of *) is not present in the scan, where pp. 18, 20 and 22 all
% carry it plainly. Same on the first note of p. 27. Set as *) like the rest.
\setcounter{footnote}{0}%
konstaterer, med mindre Iagttageren netop hviler i Forhold til Æteren.

Da \emph{Ritz}' Standpunkt\footnote{Sml. P. \emph{Ehrenfest}:
Æterhypotesens Krise. Fys. Tidsskr. XII, 242, 1913--14.} ikke kan
fastholdes, haves for Tiden kun disse to Forsøg paa en Løsning af
Vanskelighederne; \emph{Wiechert} har kaldt dem det betingede og det
ubetingede Relativitetsprincip. Da de fører til identiske iagttagelige
Resultater, maa det, som Lorentz mener, for en Fysiker blive en
Smagssag, hvilken Opfattelse han foretrækker. Einsteins Synsmaade er
fænomenologisk; dens store formelle Simpelhed har betinget, at først
gennem den overskuede man alle Relativitetsopfattelsens Konsekvenser.
Lorentz' Behandlingsmaade er af en mere atomistisk Natur; den søger at
give os en Antydning af en Forstaaelse af de paagældende Fænomeners
Mekanisme, fører os til i de negative Resultater af Michelson- og
Trouton-Noble-Forsøget og i de negative Resultater, som vi ogsaa
venter, at eventuelle nye tilsvarende Forsøg paa at paavise Jordens
Translation vil give, at se en Kompensation af modsatte Virkninger.

5.\ Her skal indskydes et Par Bemærkninger om, hvorvidt
Lorentz-Einsteins Opfattelse bekræftes af Erfaringen eller har
Mulighed for at blive det. En Afgørelse ad Erfaringens Vej af, hvilket
af de to Standpunkter der er \guillemotright
rigtigt\guillemotleft{}, er i Følge Sagens Natur udelukket.

Det kommer da først i Betragtning, at Michelson- og
Trouton-Noble-Forsøget hidtil ikke har kunnet forklares paa anden
Maade. Videre har især Undersøgelser af Elektronmassens (den
transversale Masses) Variation med Hastigheden faaet Betydning. Her har
nye Forsøg af \emph{Schaefer} og \emph{Neumann} (1913) ført til en fra
alle Sider anerkendt Bekræftelse af Relativitetsteoriens Fordringer; i
det mindste har selv \emph{Kaufmann}, der anstillede de første Forsøg
herover og længe mente, at hans Resultater var uforenelige med
Relativitetsteoriens Fordringer, nu anerkendt det fuldtud (p.~1227 i
hans nye Elektricitetslære, Müller-Pouillet, Bd.~IV). Endnu nyere
Forsøg af \emph{Guye} og \emph{Lavanchy} (1915) har givet samme
Resultat.

I denne Forbindelse er det rimeligt at erindre om, at de Maxwellske
Ligninger, som ansaas for et rigtigt Udtryk for vore Erfaringer
angaaende elektrodynamiske Fænomener i Vakuum, tilfredsstiller
Relativitetsprincipet; men paa den anden Side har jo rigtignok den
% --- p. 24 ---
\opage{24} % No footnotes on this page.
Newtonske Mekanik, som ansaas for et rigtigt Udtryk for vore
Erfaringer angaaende mekaniske Fænomener, maattet forandres.
Imidlertid er selv den efter almindelig Opfattelse uhyre Nøjagtighed,
hvormed den Newtonske Mekanik er bekræftet i Astronomien, dog ikke
større, eller rettere den Forandring i Mekaniken, som
Relativitetsprincipet forlanger, kvantitativt saa lille, at
Overensstemmelsen med den nye Mekanik bliver lige saa god, ja i
Virkeligheden bedre. Thi den gamle Mekanik formaaede som bekendt ikke
at forklare et Beløb paa $43''$ pr. Aarhundrede i Merkurs
Perihelbevægelse. Disse $43''$ kan, som nylig (1915) vist af Einstein,
forklares ved hans udvidede Relativitetsteori, der ogsaa omfatter
Gravitationen, men desværre kan der ikke gaas nærmere ind paa dette
her.

Det er vel endelig ikke overflødigt at erindre om, at der naturligvis
intet Forsøg kendes, som strider mod Relativitetsprincipet; i saa Fald
maatte dette jo, naar Forsøget var paalideligt gennemført, straks
falde; det vilde altsaa tillige betyde, at ogsaa de Lorentzske
Hypoteser maatte opgives eller ændres.

Hvad Eksperimenterne i Øjeblikket kan sige angaaende
Relativitetsprincipet, maa da sammenfattes derhen, at det foreløbigt
opstillede Relativitetsprincip ikke endnu har ført til iagttagelige
Konsekvenser, som strider mod Erfaringen, men at denne tværtimod i de
ganske vist ret faa Tilfælde, hvor en eksperimentel Prøve hidtil har
været mulig, har bekræftet det.

6.\ Jeg skal derefter gaa over til nærmere at omtale enkelte
Hovedpunkter af Professor Kromans interessante Arbejde. Hovedsagen er
naturligvis Forsøget (S.~15) paa at paavise, at Einsteins Lære strider
mod Logikens Modsigelsessætning.

Der er Tale om to Iagttagere $A$ og $B$, hver hørende til sit af to
Systemer, der bevæger sig i Forhold til hinanden med den konstante
Hastighed $v$; $A$ og $B$ har hver i sit System ordnet sine Ures
Visning under den Forudsætning, at Lysets Hastighed for hver af dem er
$c$ i alle Retninger. Efter Einstein har hver af dem Ret i sin
Fremfærd, da ingen af dem har Acceleration. Der indvendes nu: Her er
altsaa to Sandheder, men dette strider mod Logikens
Modsigelsessætning. Efter min Mening er det nu ikke uden videre
indlysende, at dette er Tilfældet; de to Ting, der skulde være
uforenelige, er to forskellige Iagttageres Fremgangsmaade over for to
forskellige Opgaver, nemlig Indstillingen af hver Iagttagers Sæt Ure.
Imidlertid skal Mod\-%
% --- p. 25 ---
\opage{25}% No footnotes on this page.
sætningen komme tydeligt frem, naar man spørger, hvad $B$ skulde have
gjort, hvis han ogsaa havde troet sit System bevæget (uden
Acceleration). Men dette Spørgsmaal er meningsløst (stadig under
Forudsætning af, at Erfaringen som hidtil vedblivende viser det umuligt
at fastslaa en absolut Bevægelse i Forhold til en eventuel Æter). Hvis
$B$ skulde have troet sig bevæget, maatte han for at kunne gøre noget
ved de foreliggende Spørgsmaal have kendt sin Hastighed, hvilket han er
ude af Stand til. Hans eneste mulige Standpunkt er at lade, som om hans
System hviler, enten han saa med Einstein mener, at der overhovedet
ikke er Mening i at tale om absolut Bevægelse, eller med Lorentz, at
han blot er ude af Stand til at faa at vide, om han hviler i Forhold
til Æteren eller ej. At $B$ skulde give $A$ Fortrin for sig selv og
antage $A$ hvilende, kan der ingen Grund anføres for.

\guillemotright $B$'s System er forkortet, siger $A$. Nej, siger
$B$.\guillemotleft{} Mellem disse to Standpunkter forlanges det, at
Fysikeren som Tredjemand skal skaffe Enhed. Men som foran berørt er der
egentlig intet at strides om mellem $A$ og $B$. Det kan ikke gøre $B$
noget, at $A$ mener om hans System, at Maalestokkene er forkortede og
Urene gaar for langsomt og viser Stedtid. $B$ har ved at indrette sig,
som han gjorde, faaet Enhed i sit System og det paa en saadan Maade, at
han finder simple og overskuelige Naturlove. Og naar $A$ og $B$ er
Fysikere (d.~v.~s. begge kender Resultatet af Michelson-Forsøget og de
øvrige tilsvarende Forsøg), ved hver af dem jo, hvorfor den anden bærer
sig ad som han gør.

Imidlertid har, som vi har set, Lorentz været den Tredjemand, som kunde
skaffe den søgte Enhed mellem $A$'s og $B$'s Standpunkt; de der
forlanger denne Enhed, stiller sig følgelig blot paa Lorentz'
Standpunkt. Kortere end det blev sammenfattet i 4., kan det maaske
siges saaledes: Lorentz' Antagelser om Elektroner, Kræfter o.~s.~v. (se
2.) fører til, at i Systemer der bevæger sig i Forhold til Æteren,
trækker alle Længder sig sammen og forandrer alle Ure deres Gang paa en
saadan Maade, at Iagttagere i to forskellige Systemer kommer til at
dømme om hinanden, som Relativitetssætningerne fordrer.

Det vil nu forstaas, at Lorentz og Einstein maa være enige om en Rejses
foryngende Virkning (Kroman S.~28). Den rejsende Mand behøver efter
Lorentz' Mening ikke at stille sit Ur til at gaa langsommere, det
besørger Bevægelsen i Forhold til Æteren gennem den
% --- p. 26 ---
\opage{26} % No footnotes on this page.
derved fremkaldte Forandring i Kræfterne af sig selv, og denne samme
Forandring af Kræfterne maa ogsaa faa de organiske Processer til at
forløbe langsommere. Einstein derimod giver i Overensstemmelse med sit
fænomenologiske Standpunkt ingen Forklaring paa, hvordan det skal
tænkes at \guillemotright gaa til\guillemotleft . Fænomenet er en
Konsekvens af Forudsætningen, af Relativitetsprincipet, og i samme Grad
som de af Relativitetsprincipets Konsekvenser, som kan prøves, viser
sig at slaa til, i samme Grad nødes man til at tro ogsaa paa denne
specielle Konsekvens, selv om denne i Øjeblikket vanskeligt kan tænkes
prøvet. Tilfredsstiller et saadant fænomenologisk Standpunkt ikke, har
man den Udvej at holde sig til Lorentz' Opfattelse eller gøre sig andre
lignende Forestillinger om Fænomenets Mekanisme.

S.~21 indvendes, at Tanken ikke kan fatte, at et Lyspunkt samtidig
skulde kunne være paa to Steder. Denne Indvending forstaar jeg næppe
helt. (At Transformationsligningerne, som er fundne ud fra den
Forudsætning, at Lysets Hastighed i begge Systemer i alle Retninger er
$c$, fører til det Resultat, at den er $c$, er jo ikke mærkeligt). Fra
Lorentz' Standpunkt er Sagen imidlertid simpel nok: i Følge dette kan
vi altid tænke os angivet entydigt, hvor i Æteren Lysblinket er naaet
hen, dets absolutte Koordinater. Blot kan hverken $A$ eller $B$ maale
andet end dets relative Koordinater, nemlig fastslaa, hvor Lysblinket
befinder sig i deres Systemer. De finder derfor hver sit Udtryk for,
hvor det er, men disse Udtryk betegner til samme absolutte Tid det
samme Sted i Æteren, det samme absolutte Sted. Einstein kender intet
Middel til at angive Lysblinkets absolutte Sted og kan som ofte nævnt
ikke forbinde nogen Mening med at tale om dets absolutte Sted. Men naar
Lysglimtet passerer $A$ (Koordinat $x$ i $A$'s System) i det Øjeblik,
Klokken dér er $t$, saa tillader Formlerne entydigt at udregne, hvor
$B$, som vi tænker os passere $A$ i samme Øjeblik, befinder sig i sit
System (Koordinat $x'$), og hvad han mener, Klokken dér er ($t'$). Og
han vil da mene, at Lysblinket netop til Tiden $t'$ er naaet til $x'$.

7.\ Endnu skal kun knyttes nogle Bemærkninger til Professor Kromans
Omtale (S.~17--19) af Relativitetsprincipet i Newtons Mekanik. Ingen
Relativist vil benægte, at Einstein har udtalt noget nyt, som ikke
findes hos Newton, men ingen Fysiker vil paa den anden Side benægte, at
den Newtonske Mekanik ogsaa har \emph{sit} Relativitetsprincip. Stiller
vi os paa et rent mekanisk Standpunkt, og tror vi
% --- p. 27 ---
\opage{27} % FOOTNOTES on this page: two, printed *) and **). The mark of the first is a
% BARE ASTERISK in the scan -- no closing paren; cf. p. 23.
\setcounter{footnote}{0}%
paa Newtons Mekanik, saa kommer vi ikke uden om, at Hvile og jævn
retlinet Bevægelse er ganske ligestillede Begreber, er \guillemotright
fysisk ét og det samme\guillemotleft . Hvad der for én Iagttager er
Hvile, er for en anden lige saa berettiget Iagttager en jævn retlinet
Bevægelse. Derimellem kan aldrig træffes Afgørelse. At dette er
rigtigt, ses let af det matematiske Udtryk for Newtons
Mekanik\footnote{Sml. S.~252 i min tidligere Artikel (Fys. Tidsskr. X,
251, 1911--12).}. Men her skal jeg nøjes med at omtale Eksemplet S.~19
nærmere.

Der er Tale om en Vogn af Massen $m$, der fra Tiden 0 til Tiden $t$
trækkes med Accelerationen $\gamma$ uden Gnidning hen ad et Skibs Dæk.
I 1.~Tilfælde hviler Skibet, i 2.~Tilfælde sejler det med konstant
Hastighed $v$ i samme Retning, som Vognen trækkes. At der ikke er
Identitet mellem de to Tilfælde, er givet nok, men det forlanger
Mekanikens Relativitetsprincip rigtignok heller ikke. Det forlanger
kun, at en medbevæget Iagttager, en Iagttager paa Skibet skal finde
Identitet mellem, hvad der sker paa Skibet i de to Tilfælde, og dette
er rigtigt. Hellere end at lade det være en flyvende Fugl, som trækker
Vognen, vil vi foreløbig lade en Mand i Forstavnen af Skibet trække
den. Han udfører (til alt Held for Sømændene) det samme Arbejde, enten
Skibet ligger stille eller bevæger sig, nemlig Kraften $m\gamma$ Gange
Vejen $\tfrac{1}{2}\gamma t^{2}$\footnote{At Vejen i begge Tilfælde er
$\tfrac{1}{2}\gamma t^{2}$ ses lettest ved at huske paa, at saa meget
maa Manden forkorte Snoren, hvori han trækker Vognen.}, og for ham og i
det hele for Iagttagere paa Skibet, altsaa medbevægede Iagttagere, er
Vognens kinetiske Energi ogsaa i begge Tilfælde
$\tfrac{1}{2}m\gamma^{2}t^{2}$; et saa stort Arbejde kan Vognen i begge
Tilfælde udføre, før den igen staar stille paa Skibet. Men for
Iagttagere udefra er der som sagt ikke Identitet; vi kan udvide
Eksemplet til at oplyse Forholdet helt ved at tænke os en Vogn til paa
Skibet, ens med den omtalte og behandlet paa samme Maade, blot trækkes
den af en Mand, der sidder paa Land, f.~Eks. paa et Bolværk, som Skibet
netop glider forbi; lad ham hedde $L$, medens Sømanden, der trak den
første Vogn, hedder $S$. Det Arbejde, $L$ maa udføre, er forskelligt i
de to Tilfælde, paa det stilleliggende Skib naturligvis som før
$\tfrac{1}{2}m\gamma^{2}t^{2}$ men paa det bevægede Skib
$(\tfrac{1}{2}\gamma t^{2}+vt)\,m\gamma$, altsaa $m\gamma vt$ mere end
før, stemmende med, at $L$ mener, at Vognens kinetiske Energi nu bliver
$\tfrac{1}{2}m(\gamma t+v)^{2}=\tfrac{1}{2}m\gamma^{2}t^{2}+m\gamma
vt+\tfrac{1}{2}mv^{2}$. Det sidste Led $\tfrac{1}{2}mv^{2}$ hidrører jo
fra, at Vognen deltager i Skibets Bevægelse med Hastigheden $v$. $L$
maa nu mene, at den anden Vogn har den
% --- p. 28 ---
\opage{28} % No footnotes on this page.
% The running head here reads »H. M Hansen:« -- the point after M is dropped,
% where pp. 20, 22, 24 and 26 all print »H. M. Hansen:«. Running heads are not
% transcribed; recorded for the record only.
samme kinetiske Energi; Spørgsmaalet er da for $L$, hvorfra
Energimængden $m\gamma vt$ for denne Vogn stammer; den stammer fra
Skibet, der har ført $S$ Vejen $vt$ frem mod Kraften $m\gamma$, som han
trak i Snoren med, og som Snoren derfor ogsaa trak i ham med. Men
omvendt forstaar $S$ maaske ikke, hvorfor $L$ har maattet slide
$m\gamma vt$ mere med sin Vogn, da denne jo for $S$ ogsaa kun har
Energien $\tfrac{1}{2}m\gamma^{2}t^{2}$. Det overskydende Arbejde er
midlertid brugt til at flytte $L$ Vejen $vt$ tilbage, set fra eller i
% PRINTER'S DEFECT, transcribed as printed: »er midlertid« for »er
% imidlertid«. The initial i of imidlertid is not there -- clean white space
% between »er« and »midlertid«, verified at 600 ppi against the same word set
% correctly on pp. 21, 22, 23, 24 and 25.
Forhold til Skibet, idet han under denne Bevægelse var paavirket af
Kraften $m\gamma$. $L$'s og $S$'s Synsmaade er naturligvis lige
berettigede, lige rigtige; sandsynligvis har ingen af dem hvilet
absolut, eller rettere, vi mangler som før nævnt i hvert fald i
Mekaniken det fra de andre udmærkede Koordinatsystem, der kunde
definere absolut Hvile og Bevægelse.

Fuglen i det oprindelige Eksempel svarer til $L$ eller $S$, alt
eftersom man regner den hørende til Omverdenen eller til Skibet. Naar
Eksemplet her er behandlet saa udførligt, er det ogsaa for at gøre
opmærksom paa, at det, at man besidder Mekanikens
Relativitetsprincip, godt kan forenes med, at man analyserer de
enkelte Tilfælde nærmere og i Enkelthederne forfølger, hvorledes det
hænger sammen med dette Relativitetsprincip; en lignende Bemærkning
kunde gøres til Redaktør Helge Holst's Omtale af Fizeaus Forsøg.

8.\ Endelig skal jeg ganske kort berøre Redaktør \emph{Helge Holst}'s
Artikel, til hvilken jo ogsaa adskillgt af det foranstaaende har
% PRINTER'S DEFECT, transcribed as printed: »adskillgt« for »adskilligt«.
% The i is simply not set; the surrounding letters are clean at 600 ppi.
Adresse. Den af Redaktør Holst hævdede Opfattelse af
Samtidighedsbegrebets Indhold falder aabenbart sammen med den af
Lorentz antydede, og der er ingen Tvivl om, at denne Opfattelse
\emph{kan} fastholdes. Men det maa bemærkes, at det afgørende jo slet
ikke er, at det at vi ikke raader over hurtigere Bud end Lysstraalen,
faar Betydning over for den \emph{Nøjagtighed}, hvormed vi kan udføre
Samtidighedsbestemmelser for Begivenheder paa forskellige Steder (H.
H. S.~3). Vanskeligheden kommer, efter at Michelson-Forsøget i Strid
med, hvad man vilde vente, har lært os, at Lysets Hastighed i alle
Retninger her paa Jorden er $c$. Thi man kan jo ikke godt tænke sig,
at Michelson-Forsøget vilde give andet Resultat f.~Eks. paa Mars end
her paa Jorden; men dermed er vi inde i den principielle
Uoverensstemmelse mellem vor og Marsmenneskernes Tidsopfattelse, igen
inde i Diskussionen mellem $A$ og $B$. Hvorvidt man da vil tillægge den
fremsatte Opfattelse af Samtidighedsbegrebets absolutte
% --- p. 29 ---
\opage{29} % Last page of the article. No footnotes. It ends with the closing paragraph
% below -- no rule, no signature, no dateline, no »Sluttet«/date line; the
% remainder of the page is blank, and printed p. 30 opens a new item,
% »Fysik i Skolen«.
Indhold Mening (eller Interesse), bliver for Fysikere som før nævnt en
Smagssag.

Den anden Fremgangsmaade, Anvendelsen af Ure, der tages om Bord i de
bevægede Systemer, forandrer intet heri. Paa den Maade tilvejebringes
en saadan Overensstemmelse mellem de forskellige Systemers Tidsenhed
(resp. ved Maalestokoverføring Overensstemmelse i Længdeenheden), at
dersom Systemerne bevæger sig med samme Hastighed, altsaa er i relativ
Hvile, maaler de med samme Enhed. Men naar de faar relativ Hastighed,
konstaterer alligevel en Iagttager i det ene System de kendte
Uoverensstemmelser mellem hans og det andet Systems Længde- og
Tidsenheder o.~s.~v., som er udtrykt ved
Lorentz-Einstein-Transformationsformlerne. (Vi forudsætter selvfølgelig
stadig, at Relativitetssætningerne gælder, d.~v.~s. at alle Forsøg af
Michelson- og Trouton-Noble-Forsøgets Art giver negative Resultater).
Den Aarsag til Forandringerne, som Lorentz søger i Ætervinden, søger
Einstein --- i Overensstemmelse med sit fænomenologiske Standpunkt ---
i Rummets og Tidens Natur, uden at han gør sig
Mekanismeforestillinger.

% The printing sets about one extra line of white here -- a break within the
% section, without rule or ornament. Set as \medskip.
\medskip

Man maa med Interesse afvente yderligere Indlæg fra erkendelsesteoretisk
Side i disse Spørgsmaal, bygget i samme Grad som de her omhandlede
Arbejder paa en omfattende Viden om Spørgsmaalets fysiske Side.
Samtidig maa Fysikernes Arbejde være rettet mod at finde nye
Konsekvenser af Relativitetsprincipet, som egner sig til at prøves
eksperimentelt, og prøve dem, for at fastslaa endnu bedre end hidtil
sket, at vi ikke fysisk kan skelne mellem Hvile og Bevægelse, eller
altsaa at Udgangspunktet, Relativitetsprincipet, er et rigtigt Udtryk
for vor fysiske Viden om disse Spørgsmaal.

Vi har her slet ikke kunnet gaa ind paa \emph{Einsteins} nyere
almindeligere Relativitetsteori, hvor accelererede Bevægelser
omhandles og Gravitationen ogsaa er indordnet, men dertil vil der
antagelig blive Lejlighed i en senere Artikel.

\end{document}
